\documentclass[a4paper,12pt]{article}

\usepackage[T2A]{fontenc}
\usepackage[utf8]{inputenc}
\usepackage[russian]{babel}

\usepackage{amsmath,amssymb,mathtools}
\usepackage{helvet}
\renewcommand{\familydefault}{\sfdefault}

\usepackage{icomma}
\usepackage[a4paper,margin=2.1cm,headheight=15pt]{geometry}
\usepackage{microtype}

\usepackage{tikz}
\usepackage{pgfplots}
\pgfplotsset{compat=1.18}

\usepackage{tabularx,booktabs}
\usepackage{enumitem}
\usepackage{hyperref}
\usepackage{parskip}
\usepackage{fancyhdr}
\usepackage{setspace}
\usepackage{xcolor}
\usepackage{titlesec}

% -------------------------------------------------
% Палитра
% -------------------------------------------------
\definecolor{accent}{HTML}{008080}
\definecolor{darkaccent}{HTML}{004D4D}
\definecolor{textgray}{HTML}{333333}
\definecolor{softgray}{HTML}{EEF3F3}
\definecolor{linegray}{HTML}{B9C8C8}

% -------------------------------------------------
% Общая типографика
% -------------------------------------------------
\color{textgray}
\onehalfspacing
\setlength{\parindent}{0pt}
\setlength{\parskip}{0.65em}

\setlist[itemize]{
  leftmargin=1.6em,
  itemsep=0.25em,
  topsep=0.3em
}

\setlist[enumerate]{
  leftmargin=1.8em,
  itemsep=0.45em,
  topsep=0.4em
}

\titleformat{\section}
  {\Large\bfseries\color{darkaccent}}
  {\thesection.}{0.55em}{}

\titleformat{\subsection}
  {\large\bfseries\color{accent}}
  {\thesubsection.}{0.5em}{}

\titlespacing*{\section}{0pt}{1.3em}{0.5em}
\titlespacing*{\subsection}{0pt}{1em}{0.35em}

% -------------------------------------------------
% Колонтитулы
% -------------------------------------------------
\pagestyle{fancy}
\fancyhf{}
\fancyhead[L]{\small\textcolor{textgray}{Графики функций с модулем}}
\fancyhead[R]{\small\textcolor{textgray}{София Кальней}}
\fancyfoot[C]{\small\thepage}
\renewcommand{\headrulewidth}{0.4pt}
\renewcommand{\headrule}{%
  \hbox to\headwidth{\color{linegray}\leaders\hrule height \headrulewidth\hfill}
}

% -------------------------------------------------
% Служебные команды
% -------------------------------------------------
\newcommand{\key}[1]{\textcolor{darkaccent}{\textbf{#1}}}
\newcommand{\R}{\mathbb{R}}

\hypersetup{
  colorlinks=true,
  linkcolor=darkaccent,
  urlcolor=accent,
  pdftitle={Чек-лист: графики функций с модулем},
  pdfauthor={Лёвин Артём Александрович}
}

\begin{document}

% =================================================
% ТИТУЛЬНЫЙ ЛИСТ
% =================================================
\begin{titlepage}
\thispagestyle{empty}

\begin{center}

\vspace*{2.4cm}

{\Large\color{accent}\textbf{ЧЕК-ЛИСТ}}

\vspace{0.9cm}

{\Huge\bfseries\color{darkaccent}
Графики функций\\[0.25em]
с модулем}

\vspace{0.8cm}

{\Large
Раскрытие модуля, кусочная запись,\\
преобразование графиков и задачи с параметром}

\vspace{2.1cm}

\rule{0.62\linewidth}{0.6pt}

\vspace{1.2cm}

{\large
\textbf{Лёвин Артём Александрович}\\[0.4em]
эксклюзивно для Софии Кальней}

\vspace{1.5cm}

{\large ОГЭ по математике}

\vfill

{\large \today}

\end{center}

\begin{tikzpicture}[remember picture,overlay]
\draw[
  accent!45,
  line width=1.3pt
]
([xshift=2.0cm,yshift=1.65cm]current page.south west)
.. controls
([xshift=6.0cm,yshift=2.45cm]current page.south west)
and
([xshift=-6.0cm,yshift=0.95cm]current page.south east)
..
([xshift=-2.0cm,yshift=1.65cm]current page.south east);
\end{tikzpicture}

\end{titlepage}

% =================================================
% ОСНОВНОЙ ЧЕК-ЛИСТ
% =================================================

\section{Главная идея модуля}

Модуль числа показывает его расстояние от нуля. Поэтому

\[
|x|=
\begin{cases}
x, & x\ge 0,\\
-x, & x<0.
\end{cases}
\]

Для произвольного выражения \(u(x)\):

\[
|u(x)|=
\begin{cases}
u(x), & u(x)\ge 0,\\
-u(x), & u(x)<0.
\end{cases}
\]

\key{Ключевой принцип:} способ раскрытия модуля определяется знаком
\emph{всего выражения, находящегося внутри модуля}.

Например,

\[
|3x+5|=
\begin{cases}
3x+5, & 3x+5\ge0,\\
-(3x+5), & 3x+5<0.
\end{cases}
\]

После решения соответствующих неравенств:

\[
|3x+5|=
\begin{cases}
3x+5, & x\ge-\dfrac53,\\[0.4em]
-3x-5, & x<-\dfrac53.
\end{cases}
\]

\subsection{Минус перед модулем}

В выражении

\[
x^2-|3x+5|
\]

знак минус относится ко всему модулю. Поэтому

\[
x^2-|3x+5|=
\begin{cases}
x^2-(3x+5), & x\ge-\dfrac53,\\[0.4em]
x^2+(3x+5), & x<-\dfrac53.
\end{cases}
\]

То есть

\[
x^2-|3x+5|=
\begin{cases}
x^2-3x-5, & x\ge-\dfrac53,\\
x^2+3x+5, & x<-\dfrac53.
\end{cases}
\]

\key{Типичная ошибка:}
учитывать знак только первого слагаемого под модулем.
При раскрытии со знаком минус в скобки помещается всё подмодульное выражение.

% -------------------------------------------------

\section{График \(y=|x|\)}

Функция

\[
y=|x|
\]

имеет кусочную запись

\[
y=
\begin{cases}
x, & x\ge0,\\
-x, & x<0.
\end{cases}
\]

Получаются две части прямых, соединённые в точке \((0;0)\).

\begin{center}
\begin{tikzpicture}
\begin{axis}[
  width=0.72\linewidth,
  height=7.3cm,
  axis lines=middle,
  unit vector ratio=1 1,
  xmin=-4.5,
  xmax=4.5,
  ymin=-1,
  ymax=4.5,
  samples=250,
  grid=both,
  minor tick num=1,
  major grid style={linegray!55},
  minor grid style={linegray!25},
  xlabel={$x$},
  ylabel={$y$},
  clip=false,
  tick label style={font=\small}
]
\addplot[
  accent,
  line width=1.4pt,
  domain=-4:4
] {abs(x)};

\addplot[
  only marks,
  mark=*,
  mark size=2.1pt,
  darkaccent
] coordinates {(0,0)};

\node[
  anchor=south west,
  text=darkaccent
] at (axis cs:1.8,2.1) {$y=|x|$};
\end{axis}
\end{tikzpicture}
\end{center}

\key{Признак:}
модуль линейного выражения часто даёт график в форме буквы \(V\).

% -------------------------------------------------

\section{Сдвиг графика}

Рассмотрим

\[
y=|x-a|.
\]

Вершина графика находится в точке

\[
(a;0).
\]

Следовательно:

\[
y=|x-2|
\]

имеет вершину \((2;0)\), поэтому график сдвинут вправо на \(2\).

Для функции

\[
y=|x+3|=|x-(-3)|
\]

вершина находится в точке \((-3;0)\), поэтому график сдвинут влево на \(3\).

\begin{center}
\begin{tikzpicture}
\begin{axis}[
  width=0.88\linewidth,
  height=7cm,
  axis lines=middle,
  unit vector ratio=1 1,
  xmin=-5.5,
  xmax=5,
  ymin=-1,
  ymax=5,
  samples=250,
  grid=both,
  minor tick num=1,
  major grid style={linegray!50},
  minor grid style={linegray!22},
  xlabel={$x$},
  ylabel={$y$},
  clip=false,
  tick label style={font=\small}
]
\addplot[
  accent,
  line width=1.35pt,
  domain=-5:5
] {abs(x-2)};

\addplot[
  darkaccent,
  dashed,
  line width=1.25pt,
  domain=-5:2
] {abs(x+3)};

\node[text=accent,anchor=south west]
  at (axis cs:2.2,0.25) {$y=|x-2|$};

\node[text=darkaccent,anchor=south east]
  at (axis cs:-3.2,0.25) {$y=|x+3|$};
\end{axis}
\end{tikzpicture}
\end{center}

\subsection{Универсальный способ}

Запоминание направления сдвига полезно, однако надёжнее найти вершину из условия

\[
x-a=0.
\]

Именно в этой точке выражение под модулем меняет знак.

% -------------------------------------------------

\section{Как строить график через кусочную запись}

Для функции

\[
y=|x-2|
\]

сначала определяем знак подмодульного выражения:

\[
x-2\ge0
\quad\Longleftrightarrow\quad
x\ge2.
\]

Получаем

\[
y=
\begin{cases}
x-2, & x\ge2,\\
2-x, & x<2.
\end{cases}
\]

Далее:

\begin{enumerate}
\item найдите точку смены знака подмодульного выражения;
\item запишите формулу каждой ветви;
\item укажите область \(x\), где действует каждая формула;
\item выберите несколько удобных значений \(x\);
\item постройте обе части на одной координатной плоскости;
\item соедините части только в тех областях, где они определены.
\end{enumerate}

\key{Важно:}
график функции каждому допустимому \(x\) сопоставляет ровно одно значение \(y\).

Если при одном и том же \(x\) построены две разные точки графика, следует проверить ограничения ветвей.

\subsection{Граничное значение}

При раскрытии модуля равенство удобно относить к одной ветви:

\[
\begin{cases}
u(x), & u(x)\ge0,\\
-u(x), & u(x)<0.
\end{cases}
\]

В самой точке \(u(x)=0\) обе формулы дают одинаковое значение \(0\).
При построении таблицы эту точку удобно использовать как точку соединения ветвей.

% -------------------------------------------------

\section{Полезное тождество}

Для любого действительного \(x\)

\[
x^2=|x|^2.
\]

Это тождество особенно полезно в дробях, содержащих одновременно \(x^2\) и \(|x|\).

Например,

\[
x^2-4|x|
=
|x|^2-4|x|
=
|x|(|x|-4).
\]

Такое преобразование позволяет увидеть общий множитель и выполнить сокращение.

% -------------------------------------------------

\section{Сокращение дроби и сохранение ограничений}

Рассмотрим функцию

\[
f(x)=
\frac{4-|x|}
{x^2-4|x|}.
\]

Сначала найдём область определения.

Знаменатель должен быть отличен от нуля:

\[
x^2-4|x|\ne0.
\]

Используем \(x^2=|x|^2\):

\[
|x|(|x|-4)\ne0.
\]

Следовательно,

\[
|x|\ne0,
\qquad
|x|\ne4.
\]

Отсюда

\[
\boxed{x\ne0,\quad x\ne-4,\quad x\ne4.}
\]

Теперь преобразуем дробь:

\[
f(x)=
\frac{4-|x|}
{|x|(|x|-4)}
=
\frac{-(|x|-4)}
{|x|(|x|-4)}
=
-\frac1{|x|}.
\]

При сокращении множителя \(|x|-4\) значения \(x=\pm4\)
\key{остаются исключёнными}.

Итог:

\[
\boxed{
f(x)=-\frac1{|x|},
\qquad
x\ne0,\ \pm4.
}
\]

\key{Правило:}
сокращение алгебраической дроби упрощает формулу, однако исходная область определения сохраняется.

% -------------------------------------------------

\section{Раскрытие модуля в обратной пропорциональности}

Для

\[
y=-\frac1{|x|}
\]

получаем две ветви.

Если \(x>0\), то \(|x|=x\):

\[
y=-\frac1x.
\]

Если \(x<0\), то \(|x|=-x\):

\[
y=-\frac1{-x}
=
\frac1x.
\]

Следовательно,

\[
y=
\begin{cases}
-\dfrac1x, & x>0,\\[0.7em]
\dfrac1x, & x<0.
\end{cases}
\]

Обе ветви расположены ниже оси \(Ox\).

У функции

\[
f(x)=\frac{4-|x|}{x^2-4|x|}
\]

дополнительно отсутствуют точки

\[
\left(-4;-\frac14\right),
\qquad
\left(4;-\frac14\right).
\]

\begin{center}
\begin{tikzpicture}
\begin{axis}[
  width=0.86\linewidth,
  height=8.2cm,
  axis lines=middle,
  unit vector ratio=1 1,
  xmin=-6.2,
  xmax=6.2,
  ymin=-4.2,
  ymax=1.2,
  samples=300,
  grid=both,
  minor tick num=1,
  major grid style={linegray!48},
  minor grid style={linegray!20},
  xlabel={$x$},
  ylabel={$y$},
  clip=false,
  tick label style={font=\small}
]

\addplot[
  accent,
  line width=1.35pt,
  domain=-6:-0.25
] {-1/abs(x)};

\addplot[
  accent,
  line width=1.35pt,
  domain=0.25:6
] {-1/abs(x)};

\addplot[
  only marks,
  mark=*,
  mark size=2.8pt,
  mark options={
    fill=white,
    draw=darkaccent,
    line width=1pt
  }
] coordinates {
  (-4,-0.25)
  (4,-0.25)
};

\node[
  text=darkaccent,
  anchor=north
] at (axis cs:4,-0.36)
{$\left(4;-\frac14\right)$};

\node[
  text=darkaccent,
  anchor=north
] at (axis cs:-4,-0.36)
{$\left(-4;-\frac14\right)$};

\end{axis}
\end{tikzpicture}
\end{center}

Точка \(x=0\) остаётся недопустимой и соответствует вертикальной асимптоте.

% -------------------------------------------------

\section{Задачи с прямой \(y=kx\)}

Формула

\[
y=kx
\]

задаёт прямую, проходящую через начало координат.

Число \(k\) является угловым коэффициентом:

\begin{itemize}
\item \(k>0\) --- прямая возрастает;
\item \(k<0\) --- прямая убывает;
\item \(k=0\) --- получаем ось \(Ox\).
\end{itemize}

Если известна точка \((x_0;y_0)\), через которую проходит такая прямая, и \(x_0\ne0\), то

\[
k=\frac{y_0}{x_0}.
\]

\subsection{Как искать значения \(k\), при которых пересечений нет}

Для задачи

\[
f(x)=
\frac{4-|x|}
{x^2-4|x|}
\]

мы получили

\[
f(x)=-\frac1{|x|},
\qquad
x\ne0,\pm4.
\]

Выколотые точки:

\[
A\left(-4;-\frac14\right),
\qquad
B\left(4;-\frac14\right).
\]

Прямая через точку \(A\):

\[
k_A=
\frac{-1/4}{-4}
=
\frac1{16}.
\]

Прямая через точку \(B\):

\[
k_B=
\frac{-1/4}{4}
=
-\frac1{16}.
\]

Эти значения являются кандидатами, поскольку соответствующее пересечение гиперболы приходится на удалённую точку.

Обязательно проверяем также

\[
k=0.
\]

При \(k=0\):

\[
y=0,
\]

а значения функции \(-1/|x|\) всегда отрицательны. Поэтому пересечений также нет.

Итак,

\[
\boxed{
k\in
\left\{
-\frac1{16},
\,0,\,
\frac1{16}
\right\}.
}
\]

\key{Важная проверка:}
прохождение прямой через выколотую точку даёт кандидата на ответ.
После этого требуется убедиться, что у этой прямой отсутствуют другие точки пересечения с графиком.

% -------------------------------------------------

\section{Алгоритм сложной графической задачи}

Если функция содержит модуль, дробь и параметр, удобно действовать по следующей схеме.

\begin{enumerate}
\item \key{Найдите исходную область определения.}

Проверьте знаменатели и остальные ограничения.

\item \key{Упростите выражение.}

Используйте тождества, разложение на множители и сокращение.

\item \key{Сохраните исходные ограничения.}

После сокращения удалённые значения \(x\) остаются недопустимыми.

\item \key{Раскройте модуль по знаку подмодульного выражения.}

Получите кусочную запись.

\item \key{Постройте каждую ветвь в её области.}

Удобно использовать характерные точки и асимптоты.

\item \key{Нанесите выколотые точки.}

Их координаты вычисляются по упрощённой формуле.

\item \key{Постройте требуемые прямые \(y=kx\).}

Учитывайте начало координат и выколотые точки.

\item \key{Вычислите кандидаты для \(k\).}

Для точки \((x_0;y_0)\):

\[
k=\frac{y_0}{x_0}.
\]

\item \key{Проверьте остальные пересечения.}

Кандидат входит в ответ только после проверки всей области определения.

\item \key{Отдельно проверьте \(k=0\).}

Этот случай легко потерять при графическом рассуждении.
\end{enumerate}

% -------------------------------------------------

\section{Типичные ошибки}

\begin{itemize}

\item Раскрытие модуля по знаку отдельного \(x\), когда под модулем находится выражение вроде \(x-2\) или \(3x+5\).

\item Потеря скобок:

\[
-|3x+5|
\]

при положительном \(3x+5\) превращается в

\[
-(3x+5).
\]

\item Использование ветви за пределами её допустимого интервала.

\item Потеря ограничений после сокращения дроби.

\item Пропуск значений \(x\), обращающих исходный знаменатель в ноль.

\item Неверная трактовка выколотой точки как обычной точки графика.

\item Автоматическое включение значения \(k\), найденного по выколотой точке, без проверки других пересечений.

\item Пропуск специального случая \(k=0\).

\item Построение обеих ветвей гиперболы по всем значениям \(x\), хотя кусочная запись задаёт разные области.

\item Игнорирование асимптот при построении эскиза обратной пропорциональности.

\end{itemize}

% =================================================
% ТРЕНИРОВОЧНЫЙ БЛОК
% =================================================
\clearpage
\section*{Тренировочный блок}

\textcolor{darkaccent}{\textbf{Задания 1--3 --- средний уровень.}}

\begin{enumerate}

\item
Запишите функцию

\[
y=|x-2|
\]

в кусочном виде и укажите координаты вершины графика.

\item
Для функции

\[
y=|x+3|
\]

укажите координаты вершины, направление сдвига относительно графика \(y=|x|\)
и значение функции при \(x=0\).

\item
Запишите функцию

\[
y=2-|x|
\]

в кусочном виде. Укажите вершину графика и множество значений функции.

\end{enumerate}

\vspace{0.5em}
\textcolor{darkaccent}{\textbf{Задания 4--9 --- уровень выше среднего.}}

\begin{enumerate}[resume]

\item
Раскройте модуль и запишите кусочную формулу:

\[
y=x^2-|3x+6|.
\]

\item
Запишите в кусочном виде функцию

\[
y=|x^2-4|.
\]

Укажите точки, в которых меняется формула.

\item
Для функции

\[
f(x)=
\frac{3-|x|}
{x^2-3|x|}
\]

найдите область определения, упростите формулу и определите координаты выколотых точек.

\item
Для функции

\[
f(x)=
\frac{3-|x|}
{x^2-3|x|}
\]

определите все значения \(k\), при которых прямая

\[
y=kx
\]

не имеет с графиком функции общих точек.

\item
Рассмотрите функцию

\[
g(x)=
\frac{|x|-2}
{x^2-2|x|}.
\]

Найдите область определения, упростите формулу и определите все значения \(k\),
при которых прямая \(y=kx\) не пересекает график функции.

\item
Для функции

\[
y=-|x+1|+2
\]

запишите кусочную формулу, найдите вершину графика и точки пересечения с осью \(Ox\).

\end{enumerate}

\vspace{0.5em}
\textcolor{darkaccent}{\textbf{Задание 10 --- сложный уровень.}}

\begin{enumerate}[resume]

\item
Дана функция

\[
q(x)=
\frac{2-|x-1|}
{(x-1)^2-2|x-1|}.
\]

\begin{enumerate}[label=\alph*)]
\item найдите область определения;
\item упростите функцию;
\item определите координаты выколотых точек;
\item найдите все значения \(k\), при которых прямая
\[
y=kx
\]
не имеет с графиком функции общих точек.
\end{enumerate}

\end{enumerate}

% =================================================
% ОТВЕТЫ
% =================================================
\clearpage
\section*{Ответы к тренировочному блоку}

\renewcommand{\arraystretch}{1.45}

\begin{table}[ht]
\centering
\caption{Краткие ответы}
\small
\begin{tabularx}{\linewidth}{@{}cX@{}}
\toprule
\textbf{\(\№\)} & \textbf{Ответ} \\
\midrule

1 &
\[
y=
\begin{cases}
x-2, & x\ge2,\\
2-x, & x<2.
\end{cases}
\qquad
\text{Вершина }(2;0).
\]
\\

2 &
Вершина \((-3;0)\); сдвиг на \(3\) единицы влево;
\[
y(0)=3.
\]
\\

3 &
\[
y=
\begin{cases}
x+2, & x<0,\\
2-x, & x\ge0.
\end{cases}
\]
Вершина \((0;2)\);
\[
y\le2.
\]
\\

4 &
\[
y=
\begin{cases}
x^2+3x+6, & x<-2,\\
x^2-3x-6, & x\ge-2.
\end{cases}
\]
\\

5 &
\[
y=
\begin{cases}
x^2-4, & x\le-2\ \text{или}\ x\ge2,\\
4-x^2, & -2<x<2.
\end{cases}
\]
Точки смены формулы:
\[
x=-2,\qquad x=2.
\]
\\

6 &
\[
D(f)=\R\setminus\{-3,0,3\},
\qquad
f(x)=-\frac1{|x|}.
\]
Выколотые точки:
\[
\left(-3;-\frac13\right),
\qquad
\left(3;-\frac13\right).
\]
\\

7 &
\[
\boxed{
k\in
\left\{
-\frac19,\,
0,\,
\frac19
\right\}.
}
\]
\\

8 &
\[
D(g)=\R\setminus\{-2,0,2\},
\qquad
g(x)=\frac1{|x|}.
\]
Выколотые точки:
\[
\left(-2;\frac12\right),
\qquad
\left(2;\frac12\right).
\]
\[
\boxed{
k\in
\left\{
-\frac14,\,
0,\,
\frac14
\right\}.
}
\]
\\

9 &
\[
y=
\begin{cases}
x+3, & x<-1,\\
1-x, & x\ge-1.
\end{cases}
\]
Вершина:
\[
(-1;2).
\]
Пересечения с \(Ox\):
\[
(-3;0),\qquad(1;0).
\]
\\

10 &
\[
D(q)=\R\setminus\{-1,1,3\}.
\]
После сокращения:
\[
q(x)=-\frac1{|x-1|}.
\]
Выколотые точки:
\[
\left(-1;-\frac12\right),
\qquad
\left(3;-\frac12\right).
\]
Точка \(x=1\) соответствует вертикальной асимптоте.
\[
\boxed{
k\in
\left\{
-\frac16,\,
0,\,
\frac12
\right\}.
}
\]
\\

\bottomrule
\end{tabularx}
\end{table}

\end{document}